\chapter{Classes and Inheritance}
\label{ch:classes}

\standfirst{How to make a class, in the image and on disk, and what the pieces
mean.}

\section{Defining a class in the image}

Select a class category in the Browser's leftmost pane and the code pane fills
with a template. Edit it and \ct{accept} it from the middle-button menu:

\begin{st}
Object subclass: #Account
    instanceVariableNames: 'balance history'
    classVariableNames: 'InterestRate'
    poolDictionaries: ''
    category: 'MyStuff-Banking'
\end{st}

That is not a declaration. It is a message---\ct|subclass:instanceVariable\-Names: classVariableNames:poolDictionaries:category:|---sent to the class \ct{Object},
which is an object, and which responds by building a new class and installing
it in \ct{Smalltalk}. Evaluating the same expression in a workspace does the
same thing. There is nothing else going on.

\begin{description}
\item[instanceVariableNames] one space-separated string. Each instance gets
  its own copy of each, initially \ct{nil}.
\item[classVariableNames] shared by the class, its metaclass, and every
  instance and subclass. Capitalised by convention because they are, from
  inside a method, name-scoped like globals.
\item[poolDictionaries] a shared namespace, mostly of historical interest.
  Leave it empty.
\item[category] which pane of the Browser it appears in. Purely
  organisational.
\end{description}

Variants of the message make the other shapes:

\begin{center}
\small
\begin{tabular}{@{}ll@{}}
\toprule
\ctp{subclass:} & ordinary, fixed instance variables\\
\ctp{variableSubclass:} & plus an indexed part of object pointers\\
\ctp{variableByteSubclass:} & plus an indexed part of bytes\\
\bottomrule
\end{tabular}
\end{center}

\ct{Array} is a \ct{variableSubclass:}, \ct{String} is a
\ct{variableByteSubclass:}. If your class holds a fixed number of named
things---which is nearly always---use the plain one.

\section{Adding methods}

Select the class, select or type a message category, and type a method into
the code pane:

\begin{st}
balance
    "Answer how much is in the account."
    ^balance
\end{st}

\ct{accept} compiles it into the running system immediately. There is no build
step and nothing to restart. If your object is already on the screen in an
Inspector, the next time you send it that message you get the new version.

A first cut at the whole class:

\begin{st}
initialize
    balance := 0.
    history := OrderedCollection new

balance
    ^balance

deposit: anAmount
    anAmount <= 0 ifTrue: [^self error: 'deposits must be positive'].
    balance := balance + anAmount.
    history add: #deposit -> anAmount

withdraw: anAmount
    anAmount > balance ifTrue: [^Error new signal: 'insufficient funds'].
    balance := balance - anAmount.
    history add: #withdraw -> anAmount

printOn: aStream
    aStream nextPutAll: 'an Account of '; print: balance
\end{st}

\section{Creating instances}

\ct{new} is a message to the class. Inherited from \ct{Behavior}, it allocates
an instance with every named variable set to \ct{nil} --- and, in the 1983
library, does \emph{not} call \ct{initialize}.

\begin{stwarn}[\ctp{new} and \ctp{initialize}]
Pharo's \ct{Object class>>new} sends \ct{initialize}. The 1983
\ct{Behavior>>new} is a primitive and nothing else. Loading a class from
\ct{lib/} or \ct{pharo/} that defines \ct{initialize} and no class-side
\ct{new}, the loader writes \ct|^super new initialize| for you and reports
that it did. A class you type into the Browser gets no such help: write it
yourself.
\end{stwarn}

\begin{st}
Account class >> new
    ^super new initialize
\end{st}

Chapter~\ref{ch:classside} explains what \ct{Account class} means and why that
method goes somewhere different from the others.

\section{Inheritance}

\begin{st}
Account subclass: #SavingsAccount
    instanceVariableNames: 'rate'
    classVariableNames: ''
    poolDictionaries: ''
    category: 'MyStuff-Banking'
\end{st}

A subclass gets every instance variable and every method of its superclass.
Override by defining a method with the same selector; call the inherited one
with \ct{super}.

\begin{st}
SavingsAccount >> printOn: aStream
    super printOn: aStream.
    aStream nextPutAll: ' at '; print: rate
\end{st}

There is single inheritance only. For sharing behaviour across unrelated
classes there are traits, which this system flattens at load time---see
Chapter~\ref{ch:beyond}.

\subsection{Abstract classes}

Nothing marks a class abstract. The convention is a method that refuses:

\begin{st}
Shape >> area
    ^self subclassResponsibility
\end{st}

which raises an error naming the class that failed to implement it.

\section{Instance variables, class variables, and the difference}

\begin{st}
Object subclass: #Counter
    instanceVariableNames: 'count'        "one per instance"
    classVariableNames: 'Total'           "one, shared by everything"
    poolDictionaries: ''
    category: 'MyStuff'
\end{st}

\ct{count} is private to each \ct{Counter}. \ct{Total} is one variable shared
by the class, its metaclass, all its instances and all its subclasses.

\begin{stwarn}[Class variables and threads]
A class variable is shared mutable state, and in this system "shared" means
across native threads. \ct{Total := Total + 1} from two workers loses counts.
Chapter~\ref{ch:contract} is the rules; the short version is that lazy
initialisation of a class variable---\ct|^Default ifNil: [Default := ...]|
---is forbidden here, because two workers each build one.
\end{stwarn}

\section{The same class, on disk}

The Browser is the fastest way to explore, but code you want to keep belongs
in a file. The same class in Tonel format, as \ct{lib/MyStuff/Account.class.st}:

\begin{st}
"
An account.  The class comment goes here, at the top of the file, and the
Browser shows it under the class's `comment' button.
"
Class {
	#name : 'Account',
	#superclass : 'Object',
	#instVars : [ 'balance', 'history' ],
	#classVars : [ 'InterestRate' ],
	#category : 'MyStuff-Banking'
}

{ #category : 'initialization' }
Account >> initialize [
	balance := 0.
	history := OrderedCollection new
]

{ #category : 'accessing' }
Account >> balance [
	^balance
]
\end{st}

Chapter~\ref{ch:packages} covers packages, profiles and how a directory of
these becomes an image. The important property, which the 1983 chunk format
did not have, is that \textbf{load order does not matter}: definitions are
read in a pass of their own before any method is compiled.

\section{Finding your way around}

The Browser is the tool, and Chapter~\ref{ch:tools} is about it. From code,
the reflective protocol is all ordinary messages:

\begin{st}
Account superclass                   "Object"
Account selectors                    "a Set of the selectors it defines"
Account allInstVarNames              "including inherited"
Account subclasses
Account allSubclasses
Account category
Account comment
Account >> #balance                  "the CompiledMethod itself"
Smalltalk at: #Account               "the class, by name"
Object allSubclasses size            "how big the system is"
\end{st}
